\documentclass[10pt]{article}
% Shared LaTeX style for MIT 18.06 exam revision notes.
% Compile topic files with Tectonic, XeLaTeX, or LuaLaTeX.

\usepackage[a4paper,margin=0.72in]{geometry}
\usepackage{fontspec}
\usepackage{amsmath,amssymb,mathtools}
\usepackage{booktabs,array,tabularx}
\usepackage{enumitem}
\usepackage{titlesec}
\usepackage{fancyhdr}
\usepackage{microtype}
\usepackage{xcolor}

\setmainfont{Times New Roman}
\setsansfont{Arial}

\setlength{\parindent}{0pt}
\setlength{\parskip}{3.6pt}
\setlength{\abovedisplayskip}{6pt}
\setlength{\belowdisplayskip}{6pt}
\setlength{\abovedisplayshortskip}{4pt}
\setlength{\belowdisplayshortskip}{4pt}
\renewcommand{\arraystretch}{1.08}
\setlength{\tabcolsep}{4.5pt}

\titleformat{\section}
  {\normalfont\large\bfseries}
  {\thesection}
  {0.55em}
  {}
\titlespacing*{\section}{0pt}{10pt}{3pt}

\titleformat{\subsection}
  {\normalfont\normalsize\bfseries}
  {\thesubsection}
  {0.5em}
  {}
\titlespacing*{\subsection}{0pt}{7pt}{2pt}

\setlist[itemize]{leftmargin=1.35em,itemsep=1pt,topsep=2pt,parsep=0pt}
\setlist[enumerate]{leftmargin=1.55em,itemsep=1pt,topsep=2pt,parsep=0pt}

\pagestyle{fancy}
\fancyhf{}
\fancyfoot[C]{\scriptsize\color{black!45}MIT 18.06 Linear Algebra Exam Notes --- \thepage}
\renewcommand{\headrulewidth}{0pt}
\renewcommand{\footrulewidth}{0pt}

\newcommand{\notesubtitle}[1]{%
  {\centering\itshape #1\par}
  \vspace{2pt}
}

\newcommand{\source}[1]{%
  {\centering\small #1\par}
  \vspace{6pt}
}

\newcommand{\labelnote}[2]{%
  \par\smallskip
  \noindent\textbf{#1.}\quad #2
  \par\smallskip
}

\newcommand{\tightlabel}[1]{%
  \par\smallskip
  \noindent\textbf{#1.}\quad
}

\newcolumntype{Y}{>{\raggedright\arraybackslash}X}
\newcolumntype{L}[1]{>{\raggedright\arraybackslash}p{#1}}

\newenvironment{notetable}
  {\small\begin{tabularx}{\linewidth}}
  {\end{tabularx}}


\begin{document}

\begin{center}
{\LARGE Eigenvalue Patterns for Common Matrix Types\par}
\end{center}
\notesubtitle{Concise exam revision table for recognizing eigenvalue behavior quickly.}
\source{Based on handwritten MIT 18.06 revision notes.}

\section{Core Idea}

\labelnote{Core idea}{Special matrix structure often predicts eigenvalue behavior before calculation. Use structure first, then verify by computation.}

This note is an original compact reference based on standard 18.06 facts.

\section{Common Patterns}

\begin{center}
\small
\begin{tabularx}{\linewidth}{@{}L{0.28\linewidth}Y Y@{}}
\toprule
Matrix type & Eigenvalue pattern & Exam use \\
\midrule
Symmetric \(S=S^{T}\) & Real eigenvalues. & Orthogonal eigenvectors; use \(S=Q\Lambda Q^{T}\). \\
Orthogonal \(Q^{T}Q=I\) & \(|\lambda|=1\). & Lengths are preserved. \\
Skew-symmetric \(A^{T}=-A\) & Pure imaginary or zero eigenvalues. & Diagonal entries are zero. \\
Projection \(P^{2}=P\) & \(\lambda=0\) or \(1\). & Range stays, nullspace collapses. \\
Positive definite \(S\) & All \(\lambda>0\). & Use for quadratic forms and pivots. \\
Rank-one \(uv^{T}\) & At most one nonzero eigenvalue. & Most directions collapse. \\
Triangular & Diagonal entries. & Read eigenvalues directly. \\
Similar matrices & Same eigenvalues. & Change coordinates without changing spectrum. \\
\bottomrule
\end{tabularx}
\end{center}

\section{Shift, Inverse, and Powers}

If \(Ax=\lambda x\), then:
\[
(A+cI)x=(\lambda+c)x.
\]
So shifting by \(cI\) shifts eigenvalues by \(c\).

If \(A\) is invertible, then:
\[
A^{-1}x=\frac{1}{\lambda}x.
\]
So inverse eigenvalues are reciprocals.

For powers:
\[
A^{k}x=\lambda^{k}x.
\]
This is why eigenvalues control long-term behavior.

\section{Stability Memory}

\begin{center}
\small
\begin{tabularx}{\linewidth}{@{}L{0.32\linewidth}Y@{}}
\toprule
Problem type & Stability test \\
\midrule
Repeated powers \(A^k\) & Use \(|\lambda|\). \\
Differential equation \(u'=Au\) & Use \(\operatorname{Re}\lambda\). \\
Projection & Stable because eigenvalues are \(0,1\). \\
Orthogonal iteration & Lengths do not grow because \(|\lambda|=1\). \\
\bottomrule
\end{tabularx}
\end{center}

\section{Common Mistakes}

\begin{center}
\small
\begin{tabularx}{\linewidth}{@{}L{0.40\linewidth}Y@{}}
\toprule
Mistake & Correct exam habit \\
\midrule
Using \(S=Q\Lambda Q^{T}\) without symmetry. & First check \(S=S^{T}\). \\
Assuming all real matrices have real eigenvalues. & Rotations can give complex eigenvalues. \\
Forgetting zero eigenvalues in projections. & Projection kills directions in the nullspace. \\
Using \(|\lambda|\) for \(u'=Au\). & Use \(\operatorname{Re}\lambda\) for continuous systems. \\
Assuming same eigenvalues means same eigenvectors. & Similar matrices preserve eigenvalues, not necessarily vector coordinates. \\
\bottomrule
\end{tabularx}
\end{center}

\section{Final Exam Checklist}

\tightlabel{Checklist}
\begin{enumerate}
  \item Identify special structure before calculating.
  \item Check symmetry, orthogonality, projection, or triangular form.
  \item Use the structure to predict eigenvalue behavior.
  \item Verify with trace, determinant, or \(\det(A-\lambda I)\).
  \item Use shifts, inverses, and powers to transform known eigenvalues.
  \item Use \(|\lambda|\) for powers and \(\operatorname{Re}\lambda\) for differential equations.
\end{enumerate}

\end{document}
